% Self-contained appendix fragment. Uses only the theorem environments and
% elementary macros in packing-unforced-proof.tex. It does not call \appendix,
% so the main document controls appendix numbering and placement.
\clearpage
\section{A weighted embedding for one Fourier step}
\label{uf:embedding}

This appendix proves the one-step estimate
\eqref{fprof:inflation} directly from its geometric and Fourier hypotheses.
In particular, the selected caps may depend arbitrarily on the child cube.
The proof does not restore the omitted caps before applying the good-tube
condition. The pin measure need not satisfy a Frostman bound in this step.

For a bounded rectangle \(A\) of positive area, write
\(m_A=|A|^{-1}\mathcal L^2|_A\). Dilations of cubes and rectangles below
are concentric. We use the Fourier convention
\(\widehat f(\xi)=\int e^{-2\pi i x\cdot\xi}f(x)\,dx\).

\begin{proposition}[Selected-cap weighted embedding]
\label{uf:embeddingprop}
Fix an integer \(K\geq1\), an index \(0\leq j<K\), and constants
\(C_0,B_0\geq1\). Let \(0<a\leq b\), and let \(P\) be an
axis-parallel square of side \(b\). Let \(\mathcal Q\) be a finite
subfamily of squares of an axis-parallel \(a\)-grid, each contained in
\(P\). For \(L\geq L_0(C_0,K)\), set
$$
 Q^+=L^{2j+2}Q,\qquad
 \widetilde Q=LQ^+,\qquad
 P^+=L^{2j+4}P,\qquad D=L^{4K+20}.
$$
Let \(\sigma\) be a finite positive Borel measure, and put
$$
 w_Q=\sigma(Q^+),\qquad M=\sigma(P^+).
$$
Let \(\Theta\) be a finite label set. For every \(\theta\in\Theta\)
choose a function \(f_\theta\in L^2(\R^2)\), a unit vector
\(u_\theta\), a frequency center \(\xi_\theta\), and an effective
unit direction \(v_\theta\), subject to the following conditions:
\begin{enumerate}
\item The Fourier support of \(f_\theta\) is contained in
$$
 \left\{\xi:
 |(\xi-\xi_\theta)\cdot u_\theta^\perp|\leq C_0/a,
 \quad |(\xi-\xi_\theta)\cdot u_\theta|\leq C_0/b\right\}.
$$
These supports are also contained in balls of radius \(C_0/a\)
whose pointwise overlap is at most \(B_0\).
\item The angular mismatch obeys
$$
 |u_\theta-v_\theta|\leq C_0 a/b.
$$
Directions may first be given consistent signs, since the rectangles
and the tests depend only on the unoriented lines.
\item For each child there are numbers
\(\varepsilon_{Q,\theta}\in\{0,1\}\). Let
\(T_{Q,\theta}\) be the rectangle centered at the center of \(Q\),
of width \(Da\) perpendicular to \(v_\theta\) and length \(Db\)
parallel to \(v_\theta\). For a number \(H\geq1\), assume only that
\begin{equation}\label{uf:embeddinggood}
 \varepsilon_{Q,\theta}=1
 \quad\Longrightarrow\quad
 \sigma(T_{Q,\theta}\cap P^+)
 \leq H\frac ab M.
 \tag{G}
\end{equation}
\end{enumerate}
The function \(f_\theta\) is the same for all children; only the
selection numbers depend on \(Q\). Then, for every \(m>0\),
\begin{align}
 &\sum_{Q\in\mathcal Q}w_Q
   \int_{Q^+}\left|\sum_{\theta\in\Theta}
           \varepsilon_{Q,\theta}f_\theta\right|^2\,dm_{Q^+}
 \label{uf:embeddingconclusion}\\
 &\quad\leq C L^{8j+14}HM
       \sum_{\theta\in\Theta}\int_{P^+}|f_\theta|^2\,dm_{P^+}
   +C_m M a^{-2}L^{-m}
       \sum_{\theta\in\Theta}\|f_\theta\|_2^2.\nonumber
\end{align}
Here \(C\) and \(C_m\) may depend on the fixed constants
\(C_0,B_0,K\), and \(C_m\) also on \(m\). They do not depend on
\(a,b,L,H\), the measure, or the numbers of children and labels.
\end{proposition}

\begin{proof}
We separate the geometric consequence of the selected tests, local
orthogonality, and the weighted reproducing-kernel estimate.

\emph{Geometry and the selected tube weights.}
The enlarged child cubes have overlap at most
\begin{equation}\label{uf:embeddingoverlap}
 \Lambda\leq C L^{4j+4}.
\end{equation}
Indeed, in each coordinate at most \(L^{2j+2}+2\) grid centers can
place a point in the corresponding enlarged interval. Also, for a
point in \(\widetilde Q\), its distance from the center of \(P\)
in the maximum norm is at most
\((1+L^{2j+3})b/2\). It follows that
\begin{equation}\label{uf:embeddingmargin}
 \dist\left(\bigcup_{Q\in\mathcal Q}\widetilde Q,
                    \R^2\setminus P^+\right)\geq Lb
\end{equation}
for \(L\geq4\). In particular all \(Q^+\) lie in \(P^+\), and
\begin{equation}\label{uf:embeddingtotalweight}
 \sum_Qw_Q\leq\Lambda M.
\end{equation}

Fix \(\theta\), and let \(U\) be any translate of a rectangle of
width \(a\) and length \(b\), oriented along \(u_\theta\).
We claim
\begin{equation}\label{uf:embeddingtubemass}
 \sum_{\substack{Q:\varepsilon_{Q,\theta}=1\\
                        \widetilde Q\cap U\ne\varnothing}}w_Q
 \leq\Lambda H\frac ab M.
\end{equation}
If the index set is empty this is immediate. Otherwise choose one of
its cubes, \(Q_0\). Every \(Q^+\) in that sum lies inside
\(T_{Q_0,\theta}\), as we now check.

Let \(q,q_0\) be the centers of \(Q,Q_0\), and choose
\(z\in\widetilde Q\cap U\), \(z_0\in\widetilde Q_0\cap U\).
For \(x\in Q^+\), write
$$
 x-q_0=(x-q)+(q-z)+(z-z_0)+(z_0-q_0).
$$
The first, second, and fourth terms have length at most
\(C L^{2j+3}a\). The third has components at most \(a\) and \(b\)
perpendicular and parallel to \(u_\theta\). Therefore
$$
 |(x-q_0)\cdot u_\theta^\perp|\leq C L^{2j+3}a,
 \qquad
 |(x-q_0)\cdot u_\theta|\leq C L^{2j+3}b.
$$
Changing to \(v_\theta\) adds at most
\(C C_0 L^{2j+3}a\) to the transverse component. The longitudinal
component stays bounded by \(C(1+C_0)L^{2j+3}b\).
Since \(2j+3\leq2K+1\), the prescribed value of \(D\) exceeds
twice these factors once \(L\geq L_0(C_0,K)\). This proves the
containment. All the cubes in question also lie in \(P^+\).
Their overlap bound and the selected test at \(Q_0\) now give
$$
 \sum_{\substack{Q:\varepsilon_{Q,\theta}=1\\
                        \widetilde Q\cap U\ne\varnothing}}\sigma(Q^+)
 \leq\Lambda\sigma(T_{Q_0,\theta}\cap P^+)
 \leq\Lambda H(a/b)M,
$$
as claimed. Thus the mass bound for \emph{every} dual rectangle in
\eqref{uf:embeddingtubemass} is a consequence of the given tests at
selected children, not an additional assumption.

\emph{Local orthogonality for a selected subset.}
Choose a Schwartz function \(v\) whose Fourier transform is supported
in the unit ball and whose modulus is at least one on the unit square
centered at zero. To obtain one, start with the inverse transform of
a smooth bump supported in a sufficiently small ball; its value can
be made nearly constant on that square and then multiplied by a fixed
constant. For a child centered at \(q\), put
$$
 v_Q(x)=v\left(\frac{x-q}{L^{2j+2}a}\right).
$$
The product \(v_Qf_\theta\) has Fourier support in the corresponding
frequency ball enlarged by at most \(a^{-1}\). These enlarged balls
have overlap at most a fixed multiple of \(B_0\). In fact, dilating
equal-radius balls by two increases their overlap by at most a factor
nine in the plane: balls contributing at a point have their original
balls inside the concentric ball of three times the radius, and an
integration of the original overlap bound gives this count.

For every subset \(A\subseteq\Theta\), Plancherel and pointwise
Cauchy--Schwarz in frequency therefore give
\begin{equation}\label{uf:embeddingsubset}
 \left\|v_Q\sum_{\theta\in A}f_\theta\right\|_2^2
 \leq C\sum_{\theta\in A}\|v_Qf_\theta\|_2^2.
\end{equation}
The constant is independent of the subset. Apply this to
\(A=\{\theta:\varepsilon_{Q,\theta}=1\}\). Restrict the left side
to \(Q^+\), and truncate the right side at \(\widetilde Q=LQ^+\).
Schwartz decay and normalization of the cube averages give, for every
\(p>0\),
\begin{align}\label{uf:embeddinglocalorthog}
 &\int_{Q^+}\left|\sum_\theta
               \varepsilon_{Q,\theta}f_\theta\right|^2\,dm_{Q^+}
 \\
 &\quad\leq C L^2
        \sum_{\theta:\varepsilon_{Q,\theta}=1}
             \int_{\widetilde Q}|f_\theta|^2\,dm_{\widetilde Q}
       +C_p a^{-2}L^{-p}\sum_\theta\|f_\theta\|_2^2.\nonumber
\end{align}
For the first term the area ratio is exactly \(L^2\). For the second,
\(|Q^+|^{-1}\leq a^{-2}\), and the square of the cutoff is bounded
by \(C_pL^{-p}\) outside \(\widetilde Q\). No unselected label
has been introduced into the main term.

\emph{The anisotropic weighted embedding.}
For a fixed label define the finite positive averaging measure
$$
 \tau_\theta=
 \sum_{Q:\varepsilon_{Q,\theta}=1}w_Qm_{\widetilde Q}.
$$
Equations \eqref{uf:embeddingtubemass} and
\eqref{uf:embeddingtotalweight} imply
\begin{equation}\label{uf:embeddingtaumass}
 \tau_\theta(U)\leq\Lambda H(a/b)M
 \quad\text{for every dual rectangle }U,
 \qquad \tau_\theta(\R^2)\leq\Lambda M.
\end{equation}

Take a smooth compact frequency bump equal to one on
\([-C_0,C_0]^2\). Translating, rotating, and rescaling it produces
a reproducing kernel \(k_\theta\) for the Fourier rectangle in the
proposition. Thus \(f_\theta=f_\theta*k_\theta\), and
\begin{gather}\label{uf:embeddingkernel}
 \|k_\theta\|_1\leq C,\\
 |k_\theta(z)|\leq\frac{C_p}{ab}
 \left(1+\frac{|z\cdot u_\theta^\perp|}{a}
          +\frac{|z\cdot u_\theta|}{b}\right)^{-p}
 \qquad(p>0).\nonumber
\end{gather}
The translation of the frequency rectangle merely modulates the kernel
and does not change these absolute-value estimates. Each \(L^2\)
function with the specified bounded Fourier support has a continuous
representative, since its Fourier transform is also integrable.
The reproducing identity and the following pointwise estimate are
understood for that representative. Cauchy--Schwarz yields
$$
 |f_\theta(x)|^2
 \leq C\int|f_\theta(y)|^2|k_\theta(x-y)|\,dy.
$$

For each fixed \(y\), partition the plane into translates of dual
rectangles \(a\)-by-\(b\) centered on the corresponding rectangular
lattice based at \(y\). The mass estimate in
\eqref{uf:embeddingtaumass}, the kernel bound with any \(p>2\), and
summation of the lattice decay give
\begin{equation}\label{uf:embeddingkerneldensity}
 \int|k_\theta(x-y)|\,d\tau_\theta(x)
 \leq\frac C{ab}\Lambda H\frac ab M
 = C\frac{\Lambda HM}{b^2}.
\end{equation}
This computation displays the single occurrence of the tube threshold.

If \(y\notin P^+\), the support margin
\eqref{uf:embeddingmargin} gives \(|x-y|\geq Lb\) on the support
of \(\tau_\theta\). Since \(a\leq b\), the anisotropic distance
in \eqref{uf:embeddingkernel} is then at least a fixed multiple of
\(L\). Its arbitrary-order decay and the total mass in
\eqref{uf:embeddingtaumass} give the additional bound
$$
 \int|k_\theta(x-y)|\,d\tau_\theta(x)
 \leq C_p\frac{\Lambda M}{ab}L^{-p}
 \qquad(y\notin P^+).
$$
Integrating the pointwise estimate for \(f_\theta\), and using
Tonelli's theorem, now proves
\begin{equation}\label{uf:embeddingweighted}
 \int|f_\theta|^2\,d\tau_\theta
 \leq C\frac{\Lambda HM}{b^2}
              \int_{P^+}|f_\theta|^2\,dy
      +C_p\frac{\Lambda M}{ab}L^{-p}\|f_\theta\|_2^2.
\end{equation}

Finally sum \eqref{uf:embeddinglocalorthog} with weights \(w_Q\),
and apply \eqref{uf:embeddingweighted} to each label. The main
coefficient, after converting to the normalized parent average, is
$$
 C L^2\Lambda H M\frac{|P^+|}{b^2}
 \leq C L^{2+4j+4+4j+8}HM
 =C L^{8j+14}HM.
$$
The local-orthogonality remainder is bounded by
\(C_p\Lambda M a^{-2}L^{-p}\sum_\theta\|f_\theta\|_2^2\).
The weighted-kernel remainder has the same bound with an additional
\(L^2\), since \((ab)^{-1}\leq a^{-2}\).
Choosing the arbitrary decay order \(p\) larger than
\(m+4j+6\), and absorbing \eqref{uf:embeddingoverlap}, gives the
stated remainder. This proves the proposition.
\end{proof}

\subsection*{Verification of the hypotheses in the unforced transfer}

Fix a spatial parent \(P=Q_{j+1}\) and one parent angular label in
Section~\ref{uf:transfer}. Apply Proposition~\ref{uf:embeddingprop}
with
$$
 a=a_j=r_j/R,\qquad b=a_{j+1}=r_{j+1}/R,\qquad
 L=R^h,\qquad \sigma=\sigma_i,\qquad H=H_{i,j}.
$$
The spatial dilations and the test dilation then coincide exactly with
those specified there, with \(K=K_0\). In particular the proposition
does not require replacing the current tests by tests on additional
parents or by a new measure.

Take
$$
 f_\theta=F_{P,\theta},\qquad
 \varepsilon_{Q,\theta}=
    \mathbf 1_{\{\theta\text{ is good for }Q\}}.
$$
The inherited identity \eqref{fprof:identity} says that each child's
function is exactly the selected sum in the proposition. For this
fixed parent the functions \(F_{P,\theta}\) do not depend on the
child. The selections are constants on the averaging cubes, even
though different cubes may choose completely different cap subsets.

The actual Fourier direction \(u_\theta\) is the direction of the
auxiliary angular interval, and \(v_\theta\) is its effective marking
direction. Their difference is \(O(\overline\delta_j)\), with
\(\overline\delta_j\leq a_j/a_{j+1}\) by
\eqref{uf:angles}. This supplies the angular mismatch hypothesis.
The marking rule supplies \eqref{uf:embeddinggood}: multiplying its
conditional probability bound by \(\sigma_i(P^+)\) gives precisely
the unnormalized inequality there. Thus the stronger-looking bound
\eqref{uf:embeddingtubemass} for arbitrary dual rectangles has been
derived from the existing marks, rather than silently assumed.

The circular support of each auxiliary cap has tangential frequency
width \(O_T(R/r_j)=O_T(a^{-1})\) and normal width
\(O_T(R/r_j^2)\leq O_T(R/r_{j+1})=O_T(b^{-1})\).
The fixed frequency convolution adds bounded width, absorbed because
\(b\leq1\). The supports lie in bounded-overlap balls of radius
\(O_T(a^{-1})\). When labels are pairs \((S,I)\), only boundedly
many standard labels \(S\) are nonzero at a direction, and the
dyadic intervals have bounded overlap. Their enclosing balls have
uniformly bounded overlap as well: consecutive angular centers at a
given level are separated by a fixed multiple of \(\delta_j\),
up to this bounded multiplicity. The constants \(C_0,B_0\) may
depend on the fixed annular ratio, hence on \(T\), but not on \(R\).
The terminal widths are at least a constant times \(R^{-1}\), so
the fixed frequency smoothing causes no growing ball multiplicity.

All functions in this application belong to \(L^2\). For example,
a bounded-density finite circular measure convolved with the fixed
compact smooth frequency function has a smooth compactly supported
Fourier representative. This observation applies also to the hard
angular indicator pieces; no differentiation of those indicators is
needed anywhere in Proposition~\ref{uf:embeddingprop}.

Consequently \eqref{uf:embeddingconclusion}, followed by summation
over parents and parent angular labels, proves
\eqref{fprof:inflation}, with a loss \(R^{C_{K_0}h}\) and the single
factor \(H_{i,j}\). The remainder is rapidly decreasing in the
required sense. Indeed \(a_j\geq R^{-1}\), and for any target
decay power \(A\), the remainder in the proposition contains
$$
 a_j^{-2}L^{-m}\leq R^{2-hm}.
$$
For fixed \(h>0\), choose \(m\) sufficiently large to absorb \(A\)
and the polynomial bounds on the remaining sums and global norms.
Constants depending on this choice are independent of the frequency.

\begin{remark}[Scope and necessary distinctions]
The proposition permits arbitrary finite positive pin measures, even
atomic ones, and arbitrary child-dependent subsets of the fixed cap
family. It would not justify allowing the selection to vary pointwise
inside a child average, substituting unrelated numbers for the masses
\(\sigma(Q^+)\), or replacing the common parent functions by
child-dependent functions. The present transfer makes none of those
changes. This appendix verifies the one-step analytic implication; it
does not establish optimality of the profile bound or of the resulting
dimension condition.
\end{remark}
